% !TeX root = ../main.tex
% ============================================================================
%  Section 1 -- Introduction
% ============================================================================

\section{Introduction}

\begin{frame}{Motivation}
  \begin{itemize}
    \item First point, resting on earlier work \cite{example2024}
    \item Second point
      \begin{itemize}
        \item A supporting detail
      \end{itemize}
    \item Third point, and two sources at once \cite{example2024,sample2023}
  \end{itemize}

  % The two \cite commands above put [1] and [2] in the text, the matching
  % full references at the bottom of this slide, and both works on the
  % references slide at the end. Note that citing [1] twice on one slide
  % still gives one line at the bottom. Delete them with the placeholder text.
\end{frame}

\begin{frame}{Contributions}
  \begin{block}{What this work does}
    A short statement of the contribution.
  \end{block}
\end{frame}

% ============================================================================
%  DELETE THIS SLIDE -- it is here so a fresh copy proves Thai is working
%  ลบสไลด์นี้ทิ้งได้ มีไว้เพื่อยืนยันว่าภาษาไทยใช้งานได้ตั้งแต่เริ่ม
% ============================================================================
\begin{frame}{Thai check / ตรวจสอบภาษาไทย}
  ถ้าคุณอ่านบรรทัดนี้ออก แปลว่าฟอนต์ภาษาไทยทำงานถูกต้องแล้ว
  พิมพ์ภาษาไทยได้ทุกที่ ทั้งหัวเรื่องและเนื้อหา โดยไม่ต้องครอบคำสั่งใด ๆ

  \vspace{0.4em}
  If you can read the line above, Thai is working. Type it anywhere ---
  titles included --- with no wrapping macro. Bold carries across:
  \textbf{Results / ผลลัพธ์}.
\end{frame}
